% ID: SEL-D0007-BA-E02
% ESTADO: validada. S008, 2026-11-07; tres pasadas del mismo autor y revisión matemática, didáctica y visual.
\textbf{a) Determinante con parámetro.}\par\nopagebreak[4]
Una matriz cuadrada tiene inversa si y sólo si su determinante no es cero. Desarrollamos por la primera fila:
\[
\det A=1(2+3)-(-1)(3m)+m(-2m)
=5+3m-2m^2=-(m+1)(2m-5).
\]
Por tanto, \boxed{A\text{ es invertible si }m\ne-1\text{ y }m\ne5/2}. En los dos valores excluidos el determinante es cero.

\textbf{b) Ecuación para m=2.}\par\nopagebreak[4]
Para $m=2$ el determinante es 3, de modo que existe inversa. La ecuación $XA-A^2=I_3$ se transforma en $XA=I_3+A^2$. Multiplicamos \emph{por la derecha}:
\[
X=(I_3+A^2)A^{-1}=A^{-1}+A.
\]
La simplificación usa asociatividad: $A^2A^{-1}=A(AA^{-1})=A$, sin intercambiar factores arbitrariamente. Para este valor,
\[
A=\begin{pmatrix}1&-1&2\\0&2&-3\\2&1&1\end{pmatrix},\qquad
A^{-1}=\begin{pmatrix}5/3&1&-1/3\\-2&-1&1\\-4/3&-1&2/3\end{pmatrix}.
\]
La inversa puede calcularse mediante Gauss-Jordan sobre $(A\mid I_3)$; sus productos con $A$ por ambos lados dan $I_3$. Sumando,
\[
\boxed{X=\begin{pmatrix}8/3&0&5/3\\-2&1&-2\\2/3&0&5/3\end{pmatrix}.}
\]
Comprobación directa de la estructura: $(A+A^{-1})A-A^2=AA+A^{-1}A-A^2=I_3$. La solución es única porque $A$ es invertible.
